\clearpage
\section{Kode Konfigurasi Benchmark}

Berikut adalah potongan kode konfigurasi dan skrip \textit{workload} yang digunakan dalam pengujian benchmark menggunakan Hyperledger Caliper.

\subsection{Konfigurasi Skenario (scenarios.json)}
Potongan konfigurasi ini menunjukkan definisi skenario untuk pengujian \textit{Throughput Step} (peningkatan beban bertahap) dan \textit{Fault Injection}.

\begin{lstlisting}[language=json, caption={Definisi Skenario Throughput dan Fault Injection}, label={lst:scenarios_json}]
{
  "scenarios": {
    "throughput-step": {
      "description": "Skenario throughput-step meningkatkan laju transaksi mint secara bertahap (5-100 TPS) guna menemukan titik jenuh pemrosesan.",
      "rounds": [
        { "id": "A1", "tps": 5 },
        { "id": "A2", "tps": 10 },
        { "id": "A3", "tps": 15 },
        { "id": "A4", "tps": 20 },
        { "id": "A5", "tps": 25 },
        { "id": "A6", "tps": 30 },
        { "id": "A7", "tps": 35 },
        { "id": "A8", "tps": 40 },
        { "id": "A9", "tps": 45 },
        { "id": "A10", "tps": 50 },
        { "id": "A11", "tps": 55 },
        { "id": "A12", "tps": 60 },
        { "id": "A13", "tps": 65 },
        { "id": "A14", "tps": 70 }
      ],
      "commonConfig": {
        "txDuration": 60,
        "rateControllerType": "fixed-rate",
        "workloadModule": "benchmarks/workload/mint-certificate-workload.js"
      }
    },
    "fault-injection": {
      "description": "Skenario fault-injection mengevaluasi ketahanan jaringan ketika salah satu validator/pod diberi gangguan terkontrol.",
      "rounds": [
        {
          "id": "F1",
          "label": "fault-tolerance",
          "rateTps": "OPTIMAL",
          "txDuration": 420,
          "workload": {
            "module": "benchmarks/workload/mint-certificate-workload.js",
            "arguments": {
              "batchId": "FAULT",
              "codeUniv": "FAULT"
            }
          }
        }
      ],
      "commonConfig": { "workers": 3 }
    }
  }
}
\end{lstlisting}

\subsection{Workload Minting (mint-certificate-workload.js)}
Skrip ini mengatur logika pembuatan transaksi penulisan (\textit{write}) untuk pencetakan sertifikat NFT dengan ID unik.

\begin{lstlisting}[language=JavaScript, caption={Workload Module untuk Minting Sertifikat}, label={lst:mint_workload}]
"use strict";

const { WorkloadModuleBase } = require("@hyperledger/caliper-core");
const { resolveAccountInfos } = require("./account-utils");

class MintSertifikatWorkload extends WorkloadModuleBase {
  constructor() {
    super();
    this.txIndex = 0;
    this.accountInfos = [];
    this.selectedAccountInfo = undefined;
  }

  async initializeWorkloadModule(workerIndex, totalWorkers, roundIndex, roundArguments, sutAdapter, sutContext) {
    await super.initializeWorkloadModule(workerIndex, totalWorkers, roundIndex, roundArguments, sutAdapter, sutContext);
    this.accountInfos = await resolveAccountInfos(this.sutAdapter.web3, totalWorkers);
    this.selectedAccountInfo = this.accountInfos[this.workerIndex % this.accountInfos.length];
  }

  async submitTransaction() {
    this.txIndex++;
    const recipientAddress = this.selectedAccountInfo.address;
    
    // Create a unique ID for each new token
    const uniqueId = this.workerIndex * 1000000 + this.txIndex;

    const contractName = "MintCertificate";
    const myArgs = {
      contract: contractName,
      contractId: contractName,
      verb: "benchmarkMint",
      args: [recipientAddress, uniqueId],
      readOnly: false,
      from: this.selectedAccountInfo.address,
    };

    await this.sutAdapter.sendRequests(myArgs);
  }
}

function createWorkloadModule() {
  return new MintSertifikatWorkload();
}

module.exports.createWorkloadModule = createWorkloadModule;
\end{lstlisting}

\subsection{Workload Reading (read-certificate-workload.js)}
Skrip ini mengatur logika pembacaan (\textit{read}) data sertifikat secara acak dari blockchain.

\begin{lstlisting}[language=JavaScript, caption={Workload Module untuk Reading Sertifikat}, label={lst:read_workload}]
"use strict";

const { WorkloadModuleBase } = require("@hyperledger/caliper-core");

class ReadSertifikatWorkload extends WorkloadModuleBase {
  constructor() { super(); }

  async initializeWorkloadModule(workerIndex, totalWorkers, roundIndex, roundArguments, sutAdapter, sutContext) {
    await super.initializeWorkloadModule(workerIndex, totalWorkers, roundIndex, roundArguments, sutAdapter, sutContext);
    
    if (!this.roundArguments.totalTokens || this.roundArguments.totalTokens === 0) {
      throw new Error("Read workload requires 'totalTokens' argument > 0.");
    }
  }

  async submitTransaction() {
    const offset = Number(this.roundArguments.tokenIdOffset || 1);
    let minId = offset;
    let maxId = offset + Number(this.roundArguments.totalTokens) - 1;

    // Generate a random token ID within [minId, maxId]
    const randomTokenId = Math.floor(Math.random() * (maxId - minId + 1)) + minId;

    const contractName = "MintCertificate";
    const myArgs = {
      contract: contractName,
      contractId: contractName,
      verb: "getCertificateNumber",
      args: [randomTokenId],
      readOnly: true,
    };

    await this.sutAdapter.sendRequests(myArgs);
  }
}

function createWorkloadModule() {
  return new ReadSertifikatWorkload();
}

module.exports.createWorkloadModule = createWorkloadModule;
\end{lstlisting}
